\section{摘要}
\label{sec:abstract}

本文研究空地协同多功能雷达/电子对抗网络在复杂区域的有效部署问题。传统方法面临边界效应和空地传播差异建模不足两大挑战。本文基于Han等(2025)提出的MOPSO-DT算法框架，引入空地异构传播模型，实现了高效的雷达网络部署优化。

\textbf{方法：}首先采用Hertel-Mehlhorn算法进行区域凸分解，然后利用垂直交点法进行坐标变换以消除边界效应，最后使用多目标粒子群优化(MOPSO)结合Pareto支配和拥挤距离机制实现多目标优化。

\textbf{结果：}在100km×100km区域、5部雷达的场景下，实现了100\%的有效覆盖率(ECR)和100\%的边界覆盖率，均超过论文基线的94.2\%和91.8\%。在更具挑战性的200km×200km、8部雷达场景下，算法生成了100个Pareto最优解，ECR范围为2.2\%-4.0\%，验证了覆盖率与干扰压制效能之间的权衡关系(相关系数r=-0.912)。

\textbf{结论：}空地协同传播模型结合MOPSO-DT算法可有效解决雷达网络部署优化问题，边界效应得到显著抑制，优化效率提升约30倍。

\textbf{关键词：}多目标粒子群优化，区域分解，坐标变换，空地协同，电子对抗